% id: PHY-0001
% title: ปริมาตรน้ำที่เหลือในถังทรงกระบอกเอียง
% camp: อื่น ๆ
% topic: ของไหล
% year: 2551
% tags: สอวน., ม.5, ทรงกระบอก, ปริมาตร
% origin: สอวน. ม.5 2551 · ข้อ 1 ในภาพที่ผู้ใช้ส่ง
% review: ถอดจากภาพแล้ว · รอตรวจทาน
% diagram: PHY-0001.svg
% scan: sources/PHY-0001-0002-original.png
\documentclass[12pt]{article}
\usepackage{amsmath,amssymb,fontspec,polyglossia,tikz}
\setdefaultlanguage{thai}
\setmainfont{Noto Serif Thai}
\usetikzlibrary{arrows.meta}
\begin{document}
\textbf{ข้อ 1.} (สอวน. ม.5 2551)
ถังน้ำทรงกระบอกรัศมี $R$ ความสูง $H$ ใส่น้ำจนเต็มเมื่อตั้งดิ่ง
จากนั้นจึงจับเอียงทำมุม $\theta$ กับแนวดิ่ง จะเหลือน้ำในถังเป็นปริมาตรเท่าใด

% รูปวาดใหม่ตามภาพต้นฉบับ ไม่ได้ระบุสเกลหรือความสูงของน้ำเพิ่ม
\begin{tikzpicture}[x=1cm,y=1cm,>=Stealth]
  \draw (-1.3,0) -- (4.5,0);
  \node[above] at (3.7,0) {พื้นระดับ};
  \filldraw[fill=gray!20] (-0.4,0) rectangle (0.25,0.5);
  \draw[dashed] (0.25,0.5) -- (0.25,2.7);
  \draw[->] (-1.05,2.35) -- (-1.05,1.65) node[right,midway] {$g$};
  \begin{scope}[shift={(0.85,0.25)},rotate=-20]
    \draw (0,0) ellipse (0.65 and 0.21);
    \draw (-0.65,0) -- (-0.65,3.1);
    \draw (0.65,0) -- (0.65,3.1);
    \draw (0,3.1) ellipse (0.65 and 0.21);
    \draw[<->] (0.93,0) -- (0.93,3.1) node[right,midway] {$H$};
    \draw[->] (0,3.1) -- (0.65,3.1) node[above right] {$R$};
  \end{scope}
  \draw (0.25,1.3) arc[start angle=90,end angle=70,radius=0.8];
  \node at (0.43,1.55) {$\theta$};
\end{tikzpicture}
\end{document}
